\documentclass[11pt,letterpaper]{article}
\usepackage[T1]{fontenc}
\usepackage{lmodern}
\usepackage[hmargin=1in,vmargin=0.95in]{geometry}
\usepackage{amsmath,amssymb,amsthm}
\usepackage{microtype}
\usepackage[hidelinks]{hyperref}

\newtheorem{theorem}{Theorem}
\hypersetup{pdftitle={A proof of Fink's conjecture on recursive divisors}}
\begin{document}
\begin{center}
{\Large\bfseries A proof of Fink's conjecture on recursive divisors\par}
\medskip
Aafi Khan Malek\\
Independent researcher\\
\href{mailto:aafimalek2032@gmail.com}{\texttt{aafimalek2032@gmail.com}}\par
\smallskip
26 November 2026
\end{center}

\begin{abstract}
We prove Conjecture 1 of T.~M.~A.~Fink's 2023 preprint on ordered
factorizations and recursive divisors. The conjectured binomial sum also
admits a nonnegative refinement that counts strict divisor chains by length.
We prove the refinement by direct counting and by formal power series.
The power series proof is formalized in Lean~4 with mathlib.
\end{abstract}

\noindent\textit{2020 Mathematics Subject Classification:} 11A51, 05A15, 06A07.\\
\textit{Keywords:} ordered factorization, recursive divisor, divisor chain,
Kalm\'ar's problem, formal verification.

\section{The identity}

A strict divisor chain from $n$ is a sequence $n=d_0>d_1>\cdots>d_\ell$
with $d_j\mid d_{j-1}$ at each step. Let $C_n(u)$ be the polynomial whose
coefficient of $u^\ell$ counts these chains, including the zero-step chain.
Choosing the first step gives
\begin{equation}
C_n(u)=1+u\sum_{\substack{d\mid n\\d<n}} C_d(u).
\label{eq:recurrence}
\end{equation}
At $u=1$ this is the defining recurrence for Fink's $\kappa_0(n)$.

\begin{theorem}
Let $p_1,\ldots,p_r$ be distinct primes and let
$n=\prod_{k=1}^r p_k^{a_k}$ with $a_k\ge 0$. Then
\begin{equation}
C_n(u)=
\sum_{i_1=0}^{a_1}\cdots\sum_{i_r=0}^{a_r}
u^{i_1+\cdots+i_r}
\prod_{k=1}^r\binom{a_k}{i_k}\binom{a_k+i_1+\cdots+i_{k-1}}{a_k}.
\label{eq:main}
\end{equation}
For $r=0$ the expression is $1$.
\end{theorem}

Setting $u=1$ gives the symmetric form of Fink's conjecture
\cite[\S\,5]{fink}. Summation over $i_r$ contributes $2^{a_r}$ and yields
the form stated as Conjecture~1. This also gives $K(n)=C_n(1)/2$ for
$n>1$, by Fink's relation between $\kappa_0$ and the ordered factorization
count $K$ \cite[\S\,1]{fink}. At $n=1$ both counts are $1$.

\section{Proof by counting factorizations}

A chain of length $\ell$ corresponds to an ordered product
$n=b_1\cdots b_\ell d$, where $b_j=d_{j-1}/d_j>1$ and $d=d_\ell\ge 1$.
Successive division by $b_1,\ldots,b_\ell$ recovers the chain.

Fix $i_k$ as the number of quotient factors whose lowest-index prime is
$p_k$, and put $s_k=i_1+\cdots+i_k$ with $s_0=0$. There are
$\prod_{k=1}^r\binom{s_k}{i_k}$ ways to order these minima among the
$\ell=s_r$ quotient slots. The factor at step $k$ interleaves the $i_k$
copies of the new minimum with the $s_{k-1}$ previously ordered minima.

For each $p_k$, reserve one copy for each slot with minimum $p_k$.
The remaining $a_k-i_k$ copies can go into any of the $s_k$ slots with
minimum of index at most $k$, or into the terminal factor $d$. Stars and
bars gives $\binom{a_k-i_k+s_k}{a_k-i_k}$ allocations. These choices are
independent across primes and uniquely determine the factorization. Since
\[
\binom{s_k}{i_k}\binom{a_k-i_k+s_k}{a_k-i_k}
=\binom{a_k}{i_k}\binom{a_k+s_{k-1}}{a_k},
\]
each product in \eqref{eq:main} counts chains with exactly the specified
quotient minima. Summing over the counts $i_k$ gives \eqref{eq:main}.
\hfill$\square$

\section{Formal power series proof}

Fix the primes and let $F_k$ be the generating series of $C_n(u)$ over
the exponents of $p_k,\ldots,p_r$, with corresponding variables
$z_k,\ldots,z_r$. Put $F_{r+1}=1$ and $D_k=\prod_{j=k}^r(1-z_j)^{-1}$.
Choosing a proper divisor amounts to choosing a subvector of exponents
whose difference from the full vector is nonzero. Thus
\eqref{eq:recurrence} gives $F_k=D_k+u(D_k-1)F_k$. Solving this equation
in the formal series ring over $\mathbb Z[u]$ gives
\begin{equation}
F_k^{-1}=(1+u)\prod_{j=k}^r(1-z_j)-u,
\qquad
F_k=\frac{F_{k+1}}{1-z_k(1+uF_{k+1})}.
\label{eq:series}
\end{equation}
The second form follows by writing the first as
$(1-z_k)F_{k+1}^{-1}-uz_k$.

For nonnegative integers $a$ and $s$, the binomial theorem gives
\begin{align*}
[z_k^a]F_k^{s+1}
&=\binom{a+s}{a}F_{k+1}^{s+1}(1+uF_{k+1})^a\\
&=\sum_{i=0}^a \binom{a+s}{a}\binom{a}{i}u^iF_{k+1}^{s+i+1}.
\end{align*}
Starting with $F_1$ and $s=0$, extract $z_k^{a_k}$ in the order
$k=1,\ldots,r$. Each term uses $s=i_1+\cdots+i_{k-1}$ at step $k$.
Since $F_{r+1}=1$, the result is \eqref{eq:main}. \hfill$\square$

\section{Relation to prior formulas and verification}

The coefficients of $C_n$ are Fink's counts of proper-divisor iterations
\cite[\S\,4]{fink}. They also equal the associated divisor function
$c_\ell^{(1)}(n)$ of Law, Lettington and Schmidt, which counts $\ell$
ordered factors greater than $1$ followed by an arbitrary positive factor
\cite[Eq.\,(2)]{law}. Their alternating-sum formulas follow by
inclusion--exclusion. Equation~\eqref{eq:main} gives the coefficients as
sums of nonnegative products and specializes to Fink's conjectured sum.

For example, $C_{12}(u)=1+5u+7u^2+3u^3$, so $\kappa_0(12)=16$.
The accompanying Python script \texttt{verify\_divisor\_chains.py}
compares \eqref{eq:main} with \eqref{eq:recurrence} coefficient by
coefficient for every $1\le n\le 1{,}000{,}000$ using exact integer
arithmetic. All 4,626,619 checked coefficients agreed. Run the script
with Python 3.10 or later.

The accompanying Lean~4 file \texttt{FinkConjecture.lean} formalizes the
power series proof against mathlib v4.36.0. Its theorem
\texttt{divisorChains\_identity} proves \eqref{eq:main} for every weight
$u$ in a commutative ring and every finite list of distinct primes, so
$u=X$ in $\mathbb Z[X]$ gives the count by chain length. The theorem
\texttt{fink\_conjecture\_1} proves the printed form of Conjecture~1, with
the factor $2^{a_r}$. Lean reports that both proofs use only the standard
axioms \texttt{propext}, \texttt{Classical.choice} and \texttt{Quot.sound}.
The script and the Lean project are available at
\url{https://github.com/Aafimalek/novel-discoveries}.

The function $\kappa_0$ is sequence A067824 and $K$ is sequence A074206 in
the OEIS \cite{oeis}.

\section*{Use of AI tools}

The author used AI systems (OpenAI Codex and Anthropic Claude) throughout
this work. They searched the literature, proposed the least-prime grouping
behind the counting proof and the nested power series behind the second
proof, wrote the Python verification script and the Lean formalization, and
drafted the text of this note. The author checked every argument, ran the
verification script, and built the Lean proof, which the Lean kernel checks
independently of how it was produced. The author takes full responsibility
for the content.

\begin{thebibliography}{9}
\bibitem{fink} T.~M.~A. Fink,
\emph{Number of ordered factorizations and recursive divisors},
\href{https://arxiv.org/abs/2307.16691}{arXiv:2307.16691}, 2023.

\bibitem{law} A.~D. Law, M.~C. Lettington and K.~M. Schmidt,
\emph{Multifactorisations and divisor functions}, Utilitas Mathematica
\textbf{125} (2025), 43--60.
\href{https://doi.org/10.61091/um125-03}{doi:10.61091/um125-03}.

\bibitem{oeis} OEIS Foundation Inc., \emph{The On-Line Encyclopedia of
Integer Sequences}, \url{https://oeis.org}, sequences A067824 and A074206.
\end{thebibliography}
\end{document}
